\documentclass[10pt,a4paper]{amsart}
\usepackage[margin=19mm]{geometry}
\usepackage{amsmath,amssymb,mathtools}
\usepackage[scr=rsfs,cal=euler]{mathalfa}
\usepackage{xfrac}
\usepackage[unicode,hidelinks,bookmarks=false]{hyperref}
\newcommand{\ie}{i.e.\ }
\newcommand{\cf}{cf.\ }
\newcommand{\resp}{respectively\ }
\newcommand{\Subsection}{\subsection}
\providecommand{\nobreakdash}{\nobreak\hbox{-}}
\setlength{\emergencystretch}{2em}
\multlinegap0pt
\usepackage{amssymb}
\usepackage{mathtools}
\usepackage{xcolor}
\newtheorem{theorem}{Theorem}
\newtheorem{proposition}{Proposition}
\newtheorem{corollary}{Corollary}
\newtheorem{lemma}{Lemma}
\newcommand{\PGRM}{\mathrm{PGRM}}
\newcommand{\cl}{\operatorname{cl}}
\newcommand{\wt}{\mathrm{wt}}
\newcommand{\wtq}{\mathrm{wt}_q}
\title{Counterexamples to Charpin's Conjecture on BCH codes}
\author{Run Zheng, Yaoran Yang, Yutong Zhang, Haode Yan,, Maosheng Xiong}
\date{Source: arXiv:2607.28741v3 (CC BY 4.0)}
\begin{document}
\maketitle

\noindent\emph{Source and licence.} Counterexamples to Charpin's Conjecture on BCH codes. Version arXiv:2607.28741v3 (CC BY 4.0), distributed by arXiv under the Creative Commons Attribution 4.0 International licence, http://creativecommons.org/licenses/by/4.0/, as recorded in the arXiv licence metadata for this exact version and shown on the abstract page https://arxiv.org/abs/2607.28741v3. Full licence notice: \texttt{licenses/arxiv-2607.28741-CC-BY-4.0.txt}.\par\smallskip

\noindent\textit{Editorial scope.} Counted unit: the article's theorem mainth (source0.tex lines 153--162). The article's complete original proof of it is delivered with it (source0.tex lines 518--707, one \texttt{proof} environment of 190 lines, which the article places away from the statement). No mathematics is written, repaired or re-derived: the statement and the proof are the article's own lines, in the article's own notation, and no retained line is substituted or re-flowed.  Retained uncounted as the source-local context the proof needs: lines 251--253; lines 274--276; lines 323--339; lines 386--394; lines 424--436; lines 427--435; lines 497--509.  Nothing was omitted from any retained span, so the package carries no omission record.  The retained lines cite 1 works, whose entries are transcribed verbatim from the article's own bibliography source in the e-print and delivered with short editorial labels recorded in the item's reference mapping.  No retained line contains a figure environment, an image-inclusion command or drawing code, so no image is delivered and the package declares no asset dependency.  Editorial formatting only, in addition: an AMS \texttt{amsart} preamble that restates the article's own declarations and macros used by the retained lines, namely the environments \texttt{theorem}, \texttt{proposition}, \texttt{corollary}, \texttt{lemma} on separate counters, together with the standard packages those lines need.  The retained environments print in this document as Theorem 1, Proposition 1, Corollary 1, Lemma 1 in order of appearance, whereas the article prints the same items with the numbering of its own full document.  The complete native source of the article as distributed in the arXiv e-print for this exact version is bundled unchanged as \texttt{sources/206435/source0.tex}.
\begin{equation}\label{pre2}
\mathcal{C}\subseteq \mathcal{C}^{\prime} \quad \text{if and only if} \quad T(\mathcal{C}^{\prime}) \subseteq T(\mathcal{C}).
\end{equation}

\begin{equation}\label{pre1}
  T(\mathcal{C}_{(q,m,\delta)}) = \bigcup_{a=1}^{\delta-1} C_a.
\end{equation}

\begin{proposition}\label{properties}
Let $a$ and $b$ be two integers with $0\leq a,b\leq n-1$. Then the following
statements hold.
\begin{enumerate}
    \item  $a=b$ if and only if $[a]_q=[b]_q$, and $a<b$ if and
    only if $[a]_q<[b]_q$.
    \item 
    $
    [a q^k\bmod n]_q=[a]_qQ^k
$  for all $0\leq k\leq m-1$.
    \item  $a=\cl(a)$  if and only if
    $
    [a]_qQ^k\geq  [a]_q
   $ 
    for all $1\leq k\leq m-1$.
\end{enumerate}
\end{proposition}

\begin{equation}\label{eqdef}
T\left(\mathrm{PGRM}_q(\ell,m)\right)
=
\left\{
a:
1\leq a\leq n-1,\ 
\wtq(a)<m(q-1)-\ell
\right\}.
\end{equation}

\begin{corollary}\label{lemma2}
Let $\ell$ be an integer satisfying $1\leq \ell<m(q-1)$. Then every codeword
$\mathbf{c}\in\PGRM_q(\ell,m)$ satisfies
\begin{equation}\label{contra}
\wt(\mathbf{c})
\equiv 0
\pmod{q^{\left\lfloor (m-1)/\ell\right\rfloor}}
\quad\text{or}\quad
\wt(\mathbf{c})
\equiv -1
\pmod{q^{\left\lfloor (m-1)/\ell\right\rfloor}}.
\end{equation}
\end{corollary}

\begin{equation}
\wt(\mathbf{c})
\equiv 0
\pmod{q^{\left\lfloor (m-1)/\ell\right\rfloor}}
\quad\text{or}\quad
\wt(\mathbf{c})
\equiv -1
\pmod{q^{\left\lfloor (m-1)/\ell\right\rfloor}}.
\end{equation}

\begin{lemma}\textnormal{\cite{KasamiLin1972}}\label{kl}
Let $m$, $i$, and $j$  be integers with 
$
1\leq i\leq m-j-2$ and 
$0\leq j\leq m-2i.
$ Let $\delta=2^{m-1-j}-2^{m-1-j-i}-1.
$ 
Then
$\mathcal{C}_{(2,m,\delta)}$ satisfies
\begin{equation*}
d\bigl(\mathcal{C}_{(2,m,\delta)}\bigr)=\delta.
\end{equation*}
\end{lemma}

\begin{theorem}\label{mainth}
Let $q$ be a prime power, and let  $m$ be an integer with $m\ge 13$.  Set $u=\lfloor m/4\rfloor$ and
$t=\lfloor (m-1)/3\rfloor$. Let $s$ be an integer with $u\le s<t$, and let
$\delta=q^m-q^{m-1}-q^{m-1-u}-q^{s}-1$. Then
\begin{equation} \label{1:gapb}
d_B(\mathcal{C}_{(q,m,\delta)})=\delta\quad \hbox{ and }\quad 
  d(\mathcal{C}_{(q,m,\delta)}) \geq \delta+q^{s}, 
\end{equation}
with equality in the latter holding when $q = 2$.
\end{theorem} 

\begin{proof}[Proof of Theorem \ref{mainth}]
The proof is divided into four steps: we first show that $d_B(\mathcal{C}_{(q,m,\delta)})=\delta$, i.e., $\delta$ is a coset leader. Next, we show that $\mathcal{C}_{(q,m,\delta)} \subseteq
\mathrm{PGRM}_q(3,m)$. Then by combining $d(\mathcal{C}_{(q,m,\delta)}) \ge \delta$ and Corollary \ref{lemma2}, we prove \eqref{1:gapb} of Theorem \ref{mainth}. Now Finally we deal with the case that $q=2$. 

{\bf Step 1.} We first prove that $d_B(\mathcal{C}_{(q,m,\delta)})=\delta$, that is, $\delta$ is a coset leader.  
By Proposition~\ref{properties}, it suffices to prove that
\begin{equation}\label{suf}
[\delta]_qQ^k>[\delta]_q
\quad\text{for all }1\leq k\leq m-1.
\end{equation}

A direct calculation
gives the $q$-adic vector of $\delta$ as 
\begin{equation}\label{deltaq}
    [\delta]_q = \left( q-2, (q-1)^{(u-1)}, q-2, (q-1)^{(b)}, q-2, (q-1)^{(s)} \right),
\end{equation}
where $b = m - u - s - 2$ and $x^{(r)}$ denotes a block consisting of $r$ consecutive copies of
$x$.
The assumptions $u = \lfloor m/4 \rfloor$ and $m \ge 13$ imply that
\begin{equation*}
    m - 1 \le 4u + 2 < 6u.
\end{equation*}
It follows that
\begin{equation*}
    t = \left\lfloor \frac{m-1}{3} \right\rfloor \le 2u - 1.
\end{equation*}
Since $s < t$, we obtain
\begin{equation}\label{hao}
    s \le t - 1 \le 2u - 2.
\end{equation}
Consequently,
\begin{equation}\label{bgequ}
\begin{split}
    b &= m - u - s - 2 \\
    &\ge m - 3u\\ 
    &\ge u.
    \end{split}
\end{equation} 
Therefore, the first $u+1$ digits of $[\delta]_qQ^u$ and $[\delta]_q$
are
\begin{equation*}
    \left(q-2,(q-1)^{(u)}\right)\quad \hbox{and}\quad 
\left(q-2,(q-1)^{(u-1)},q-2\right), 
\end{equation*}
respectively. It follows that $[\delta]_q Q^{u} > [\delta]_q$. 
Similarly, since $s \ge u$, the first $u+1$ digits of
$[\delta]_q Q^{u+b+1}$ are also  $\left( q-2, (q-1)^{(u)} \right)$. Hence,
$
    [\delta]_q Q^{u+b+1} > [\delta]_q.
$ 


Moreover, for each integer  $k$ satisfying $1\leq k\leq u-1$,
$u+1\leq k\leq u+b$, or $u+b+2\leq k\leq m-1$, 
the first digit of $[\delta]_q Q^k$ is $q-1$, whereas the first digit of
$[\delta]_q$ is $q-2$.
This implies that
\begin{equation*}
    [\delta]_q Q^k > [\delta]_q
\end{equation*}
for all such $k$. Together with the cases $k=u$ and $k=u+b+1$ considered above, this
proves \eqref{suf} and hence establishes
$d_B(\mathcal{C}_{(q,m,\delta)})=\delta$.

{\bf Step 2.} Now we prove that
\begin{equation}\label{contain}
\mathcal{C}_{(q,m,\delta)}
\subseteq
\mathrm{PGRM}_q(3,m).
\end{equation}
By \eqref{eqdef},  the defining set of $\mathrm{PGRM}_q(3,m)$  is 
\begin{equation}\notag
T(\mathrm{PGRM}_q(3,m)) = \{ a \in \{ 1, \dots, n-1\} : \mathrm{wt}_q(a) < m(q-1)-3 \}.
\end{equation}
Let $a \in T(\mathrm{PGRM}_q(3,m))$ and let $[a]_q = (a_{m-1}, a_{m-2}, \dots, a_0)$ be its $q$-adic vector. Then we have
\begin{equation}\label{inequ}
    \sum_{i=0}^{m-1} (q - 1 - a_i) = m(q-1) - \mathrm{wt}_q(a) > 3.
\end{equation}
We now establish
\begin{equation*}
   [\cl(a)]_q < [\delta]_q.
\end{equation*}
by considering the following two cases.

First, suppose that $a_i \le q - 3$ for some integer $i$ with $0\leq i \leq  m-1$. Note that $a_i$ is precisely the first digit of $[a]_q Q^{m-1-i}$, while the first digit of $[\delta]_q$ is $q-2$. Thus, we have
\begin{equation*}
    [\cl(a)]_q \le [a]_q Q^{m-1-i} < [\delta]_q.
\end{equation*}

Now suppose that $a_i \ge q - 2$ for all integers $i$ with $0\leq i\leq m-1 $. Let
\begin{equation*}
k = |\{i : 0\leq i\leq m-1, a_i = q-2\}|
\end{equation*}
denote the number of digits of $[a]_q$  equal to $q-2$. It follows from \eqref{inequ} that $k \ge 4$.

If there exist two cyclically consecutive entries equal to $q-2$ in the $q$-adic vector $[a]_q$ of $a$, then there exists some integer $j$ with $0\leq j \leq   m-1$ such that the first two digits of $[a]_q Q^j$ are both equal to $q-2$. Observe that the first two digits of $[\delta]_q$ are $q-2$ and $q-1$, respectively.
It follows  that
\begin{equation*}
    [\mathrm{cl}(a)]_q \le [a]_q Q^j < [\delta]_q.
\end{equation*}

Otherwise, after replacing $[a]_q$ by $[a]_qQ^j$ for a suitable integer
$j$, if necessary, we may assume that the $q$-adic vector of $a$ is of the form
\begin{equation*}
    [a]_q = \left( q-2, (q-1)^{(g_1)}, q-2, (q-1)^{(g_2)}, \dots, q-2, (q-1)^{(g_k)} \right),
\end{equation*}
where $g_i \ge 1$ for all $1 \le i \le k$, and $g_1 = \min \{g_i : 1 \le i \le k\}$. It is clear that
\begin{equation*}
    \sum_{i=1}^k g_i = m - k.
\end{equation*}
Therefore, we have
\begin{equation*}
    g_1  \le \left\lfloor \frac{m-k}{k} \right\rfloor \le \left\lfloor \frac{m-4}{4} \right\rfloor = u - 1.
\end{equation*}

If $g_1<u-1$, then the first $g_1+2$ digits of $[a]_q$ and 
$[\delta]_q$ are 
\begin{equation*}
    \left(q-2,(q-1)^{(g_1)},q-2\right)\quad \hbox{and}\quad 
\left(q-2,(q-1)^{(g_1+1)}\right),
\end{equation*}
respectively. It follows that
\begin{equation} \label{hhh}  
[\mathrm{cl}(a)]_q \le [a]_q < [\delta]_q.
\end{equation}

It remains to consider the case $g_1=u-1$. Since $g_1 = \min \{g_i : 1 \le i \le k\}$, we have $g_i \ge u - 1$ for all $1 \le i \le k$. Therefore,
\begin{align*}
    g_2 &= m - k - \sum_{\substack{1\leq i\leq k\\ i\neq2}}g_i\\
       & \le m - k - (k-1)(u-1) \\
       &\le m - 1 - 3u,
\end{align*}
where the last inequality follows from $k \ge 4$. On the other hand, it follows from \eqref{hao} that
\begin{equation}\label{inequality}
    b - (m - 1 - 3u) = 2u - 1 - s \ge 1.
\end{equation}
Combining the above inequalities, we obtain $g_2 < b$. Then, by comparing the first $g_1 + g_2 + 3$ digits of $[a]_q$ and $[\delta]_q$, we obtain \eqref{hhh} again.


We have thus shown that $[\mathrm{cl}(a)]_q < [\delta]_q$ in all cases. Recalling \eqref{pre1} and $\mathrm{cl}(\delta) = \delta$, it follows that $a \in T(\mathcal{C}_{(q,m,\delta)})$. Therefore,
\begin{equation}\notag
T\bigl(\mathrm{PGRM}_q(3,m)\bigr)
\subseteq
T\bigl(\mathcal{C}_{(q,m,\delta)}\bigr).
\end{equation}
With \eqref{pre2}, this implies that \eqref{contain} holds.

{\bf Step 3.} We now proceed to prove   
$d(\mathcal{C}_{(q,m,\delta)}) \geq \delta+q^{s}$.
Suppose that $w = \mathrm{wt}(\mathbf{c})$ for some nonzero codeword $\mathbf{c} \in \mathcal{C}_{(q,m,\delta)}$. Then we first have $w \ge \delta$.
By \eqref{contain}, $\mathbf c$ is also a 
codeword of $\PGRM_q(3,m)$. Therefore,  Corollary~\ref{lemma2}
gives 
\begin{equation}\notag
    w \bmod q^t = 0 \quad \text{or} \quad w \bmod q^t = q^t - 1,
\end{equation}
where $t=\lfloor(m-1)/3\rfloor$.
 If $w \le \delta + q^s - 1$, then since $w\geq \delta$,  $t \le m - 1 - u$ and $s < t \le m - 1$, we have
\begin{equation*}
    q^t - q^s - 1 \le w \bmod q^t \le q^t - 2.
\end{equation*}
This contradicts \eqref{contra}. Therefore, we must have $w \ge \delta + q^s$, which implies that $d(\mathcal{C}_{(q,m,\delta)}) \geq \delta+q^{s}$. This proves \eqref{1:gapb} of Theorem \ref{mainth}. 

{\bf Step 4.} Finally, we prove that the lower bound on the minimum distance is attained
when $q=2$, i.e.,
\begin{equation}\label{maineq3}
d(\mathcal{C}_{(2,m,\delta)}) = \delta+{2}^s,
\end{equation}
where $\delta=2^{m-1}-2^{m-1-u}-2^s-1$. 

Since $m\geq 13$ and $u=\left\lfloor m/4\right\rfloor$, we have
$2\leq u\leq m-2$ and $m-2u\geq0$. Notice that $\delta+2^s=2^{m-1}-2^{m-1-u}-1$. 
By applying Lemma \ref{kl} with $(i,j)=(u,0)$, we obtain 
\begin{equation*}
d(\mathcal{C}_{(2,m,\delta+2^s)}) = \delta+2^s.
\end{equation*}
By \eqref{pre1},
\begin{equation*} T(\mathcal{C}_{(2,m,\delta)})  \subseteq T(\mathcal{C}_{(2,m,\delta+2^s)}), 
\end{equation*}
and hence  by \eqref{pre2}, 
\begin{equation*}  \mathcal{C}_{(2,m,\delta+2^s)} \subseteq \mathcal{C}_{(2,m,\delta)}.
\end{equation*}
It follows that 
\begin{equation}\notag
d(\mathcal{C}_{(2,m,\delta)})
\leq d(\mathcal{C}_{(2,m,\delta+2^s)})
= \delta+2^s. 
\end{equation}
Combining this with the lower bound already established yields \eqref{maineq3}.
\end{proof}

\begin{thebibliography}{99}
\bibitem[Kasami]{KasamiLin1972} T.~Kasami and S.~Lin, ``Some results on the minimum weight of primitive BCH codes (corresp.),'' {\em IEEE Trans. Inf. Theory}, vol.~18, no.~6, pp.~824--825, Nov. 1972.
\end{thebibliography}
\end{document}
